\section[一些 Fourier 理论]{一些 Fourier 理论（1 月 28 日）}
\srcpage{11}

今天，在深入 Fourier 理论之前，我们再谈一点谱分解。

我们从交换代数 $A$ 出发，造出几何对象 $\Spec A$。这些东西确切是什么，取决于你所采用的形式体系，例如你关心的是 $C^*$-代数与拓扑空间，还是交换环与仿射概形，或者别的可能性。

\srcpage{12}

上次我们看到 $A$-模 $M$ 如何在 $\Spec A$ 上「铺展开」成谱分解，从而在其上定义一个层。你可以用物理来想这件事：有一个物理系统，带观测量的代数 $A$ 与状态空间 $M$。$\Spec A$ 的定义性万有性质是：映射 $A\to\mathcal{O}(X)$ 与映射 $X\to\Spec A$ 自然一一对应。在物理中，你也许把做观测想成理解 $X$ 的几何的一种方式，而观测可能是某种到一条直线的函数（所以这里 $\Spec A$ 就是那条直线）。于是也许 $X$ 纤维化在这条直线上。直线上的观测量现在告诉我们一些关于 $X$ 的信息。

模与状态把这件事线性化：我们取 $M$ 并把它层化成层 $\ul{M}\to\Spec A$。$X$ 上的单个函数是到 $\mathbb{A}^1=\Spec k[x]$ 的映射；类似地，向量空间的单个矩阵（自同态）经由 Jordan 标准形给出 $\mathbb{A}^1$ 上的一个层。

量子力学中的观测正是这样发生的：我们没有经典力学中那样的经典相空间，只有它的线性化，即系统的状态 Hilbert 空间 $\mathcal{H}$。观测量是 $\mathcal{H}$ 上的自伴算子；谈谈 $\mathcal{H}$ 上全体算子构成的（非交换）代数 $\End(\mathcal{H})$ 也很有用。设 $\mathcal{O}$ 是一个观测量；那么，正如向量空间连同自同态那样，我们可以把 $\mathcal{O}$ 层化成 $\mathbb{R}$ 上的投影值测度（层的类比物），其中 $\mathbb{R}$ 是单个算子生成的代数的谱。

从这个视角看，状态 $|\varphi\rangle\in\mathcal{H}$ 是这个层的一个截面，即每个向量到各「纤维」上的投影，而纤维是投影值测度的像。给定一个截面，你可以问它支撑在哪里，即：你做了一次测量，它住在哪里？那就是支撑。我们还能做更精确的事：用范数。现在 $\|\varphi\|^2$ 是 $\mathbb{R}$ 上的复值测度：取 $\varphi$，投影到一个纤维上，然后取范数。适当归一化后，这是 $\mathbb{R}$ 上的概率测度。它告诉你这个测量预期住在哪里，也让你能问它的期望值是多少之类的问题，例如
\[
\int_{-\infty}^{\infty}x\,\|\operatorname{proj}_{\mathcal{H}_\lambda}|\varphi\rangle\|^2\,\mathrm{d}\lambda.
\]
这是下述事实的连续版本：在有限概率空间 $S$ 上的期望值是一个和：
\[
\frac{1}{\langle\varphi\mid\varphi\rangle}\sum_{\lambda\in S}\lambda\langle\psi_i\mid\varphi\rangle|\varphi\rangle,
\]
其中 $\{\psi_i\}$ 是 $S$-指标的特征向量基。\yizhu{原文此式如此排印，但按文意应为 $\frac{1}{\langle\varphi\mid\varphi\rangle}\sum_{\lambda\in S}\lambda\,|\langle\psi_\lambda\mid\varphi\rangle|^2$。}{A（原文 p.12）}

量子力学系统的数据中，最主要的观测量是 Hamiltonian（哈密顿量），即一个特定的自伴算子 $H$。它是能量泛函：它的本征态是系统的定态，即被时间演化固定的态；其特征值度量每个状态的能量。

「代数是观测量、模是状态空间、层是谱分解」这个视角，会在课程中反复出现。尽管我们大多只会见到交换代数（它们粘得好得多：你可以试试 $C^*$-代数，但 Gelfand–Naimark 告诉我们所有空间都是仿射的，那还粘什么？）。语境与语言会更花哨，但这套哲学仍会透出来。

\begin{center}
$\ast\quad\ast\quad\ast$
\end{center}

我们来谈谈 Fourier 理论与阿贝尔对偶，它们会带我们走过接下来几周。我们会把它看成谱分解的一个特例，其中交换算子的来源是一个阿贝尔群 $G$。

于是，固定阿贝尔群 $G$ 与表示 $T:G\to\Aut(V)\subset\End V$。我们把 $g\in G$ 在 $V$ 上的作用记作 $T_g$。通常，考虑 $G$ 在空间 $X$ 上的作用、再取 $X$ 上的函数，就能造出这类表示的好例子。此时 $G$ 通过拉回作用。正如本课程中已说过几次的，做这件事有不同层次的正则性，眼下选哪一种并不那么重要。如果这让你不安，就在需要的地方撒上「有限」这个词。

$G$ 作用的最典范的空间是它自己，这给出正则表示。我们想分解，所以得弄清谱是什么。若 $G$ 作用在一维向量空间上，我们得到一个特征标 $\chi:G\to\mathbb{C}^\times\subset\mathbb{C}$，这些就是可能出现的特征值。

令 $\widehat{G}:=\Hom_{\mathrm{Grp}}(G,\mathbb{C}^\times)$ 为全体特征标构成的集合；若你把 $G$ 做成拓扑的、代数的或别的什么，就对特征标也加上相应条件（例如连续、光滑或多项式）。事实上，暂且把 $G$ 想成有限阿贝尔群——并且尽量别去想有限阿贝尔群的分类，因为这个故事在更大的一般性下也走得通。

因为并非所有矩阵都可逆，$\End(V)$ 只是个幺半群，而 $G\to\Aut(V)\subset\End(V)$ 是幺半群映射。这结构不多；我们想要更多。具体说来，$\End(V)$ 本是一个代数，我们却把它遗忘

\srcpage{13}

成了幺半群。于是这看起来像是一个伴随的一半：
\[
\Hom_{\mathrm{Monoid}}(G,\mathrm{For}(\End(V)))=\Hom_{\mathrm{Alg}_{\mathbb{C}}}(?,\End V).
\]
把 $?$ 填上，就是幺半群上的自由代数，即群代数 $\mathbb{C}[G]$；它是由符号 $\delta_g$（$g\in G$）生成的代数，乘法为 $\delta_g\cdot\delta_h=\delta_{gh}$。所以你可以把它想成 $G$ 中元素的形式线性组合——也可以说它是函数 $G\to\mathbb{C}$ 的代数，记号 $\delta_g$ 正由此而来。这提示我们：$\mathbb{C}[G]$ 其实应当被想成函数的对偶，即测度。有限集上的典范计数测度为我们把函数与测度等同起来，但这不会推广。

对一般函数 $f_1,f_2:G\to\mathbb{C}$，乘积是卷积：
\[
(f_1*f_2)(k)=\sum_{g\in G}\left(\sum_{gh=k}f_1(g)f_2(h)\right)=\sum_{g\in G}f_1(g)f_2(kg^{-1}).
\]
如果你在分析课上学过卷积，这看起来应当很熟悉。另外，到目前为止我们并不需要 $G$ 是阿贝尔的！你还可以用推前与拉回重述这件事。

有三个映射 $G\times G\to G$：投影到第一、第二个因子，分别记作 $\pi_1$、$\pi_2$，还有乘法 $\mu$。于是可以定义外积
\[
f_1\boxtimes f_2:=\pi_1^*f_1\pi_1^*f_2:G\times G\longrightarrow\mathbb{C},
\]
而因为 $G$ 有限，我们可以沿 $\mu$ 把它推前，对纤维求和，得到通常的卷积；事实上，这正是前式告诉我们的。若你关心带有某种正则性的无限群，那种正则性可以在相应设定中为你买来推前映射。\yizhu{原文第二个因子亦印作 $\pi_1^*f_2$，按外积的定义应为 $\pi_2^*f_2$。}{A（原文 p.13）}

群代数 $\mathbb{C}[G]$ 交换，当且仅当 $G$ 阿贝尔；而 $G$ 在 $V$ 上的作用诱导出 $V$ 上的 $\mathbb{C}[G]$-作用，这多少正是由上述伴随给出的。\yizhu{原文此句两处均作「在 $G$ 上」，按上下文应为「在 $V$ 上」。}{A（原文 p.13）}关键事实是
\[
\Spec(\mathbb{C}[G],*)=\widehat{G},
\]
即 $\widehat{G}$ 上的函数就是群代数。为什么这是对的？谱的一个点是映射 $\Spec\mathbb{C}\to\Spec\mathbb{C}[G]$，展开来就是乘法意义下的映射 $\mathbb{C}[G]\to\mathbb{C}$，这给出一个特征标。

于是我们有一本字典
\[
(\mathbb{C}[G],*)\longleftrightarrow(\mathcal{O}(\widehat{G}),\cdot),
\]
左边的特征标与右边的 $\delta$-函数互换。此外，$g\in G$ 在左边通过平移作用，在右边则定义一个函数 $\widehat{g}$，而平移与乘以 $\widehat{g}$ 互换。

这是 Fourier 变换的一种形式。若 $f$ 是 $G$ 上的函数，则有反演公式
\[
f=\sum_{t\in\widehat{G}}\widehat{f}(t)\cdot\chi_t,
\]
其中 $\widehat{f}(t)$ 是 $f$ 在正交基下的系数。而 $\widehat{fg}$ 对应于 $\widehat{f}*\widehat{g}$。\yizhu{原文此句排印残缺，作「$(\widehat{f})$ corresponds to $(\widehat{f}\cdot)$」，按文意补足为逐点乘积与卷积互换。}{A（原文 p.13）}所有这些陈述都彼此等价，也都等价于那条关键陈述：群代数既是对偶上的函数（逐点乘法），又是 $G$ 上的函数（卷积）。

还有一个额外的对称性：$\widehat{G}$ 不只是一个集合，它本身就是一个阿贝尔群。\footnote{若 $G$ 有限，有限阿贝尔群的分类表明 $\widehat{G}$ 非典范地同构于 $G$。这在一般情形下不对，所以别拿它来培养直觉，除非你确实需要。}这是因为你可以逐点相乘 $G$ 上的函数：两个特征标的乘积仍是特征标。所以 $\widehat{G}$ 确实是那个对偶群。要对得起「对偶」这个名字，我们需要 $\widehat{\widehat{G}}$ 自然同构于 $G$：我们要给每个 $g\in G$ 指派一个从特征标到 $\mathbb{C}^\times$ 的函数，当然就是 $\chi\mapsto\chi(g)$。

还有更对称的说法：有两个投影
\[
\begin{tikzcd}[row sep=small]
G\times\widehat{G}\arrow[r,"\pi_2"]\arrow[d,"\pi_1"'] & \widehat{G}\\
G &
\end{tikzcd}
\]
而在 $G\times\widehat{G}$ 上有一个典范函数
\[
K(g,t):=\chi_t(g)=\chi_g(t),
\]
你可以把它想成一个乘性核。这里 $\chi_g$ 是 $g$ 给出的 $\widehat{G}$ 上的特征标。所以这个设定更对称（也许还会让你想起泛函分析课上的某些核）。现在可以把 Fourier 变换描述为拉回—推前：

\srcpage{14}

\[
\widehat{f}(t)=\pi_{2*}(\pi_1^*f\cdot\chi)(t)=\sum_{g}f(g)\overline{K(g,t)}.
\]
出现复共轭，是因为内积里有一个。这也意味着逆 Fourier 变换看起来略有不同：
\[
f(g)=\sum_{t}\widehat{f}(t)K(g,t).
\]
很好，我们已经把 $G$ 上的函数谱分解了——换种说法，我们同时对角化了 $G$ 在 $G$ 上函数空间的作用。「同时对角化」提醒我们：阿贝尔性对整个故事有多重要。

现在，为单个表示做这么多工作有点浪费；但因为所有 $G$-表示都是 $\mathbb{C}[G]$-表示，所以对任何表示 $\mathbb{C}[G]\to\End V$，我们都得到 $V$ 在谱 $\widehat{G}$ 上的谱分解，即
\[
V=\bigoplus_{t\in\widehat{G}}V_{\chi_t}.
\]
也就是说，$t$ 处的纤维是 $V$ 的 $\chi_t$-同型分量。对有限阿贝尔群来说这有点杀鸡用牛刀，但它可以推广。

设 $G$ 为局部紧拓扑阿贝尔群（LCA）。有许多不是有限的好例子：$\mathbb{Z}$、$U_1$、$\mathbb{R}$、$\mathbb{R}^n$、$\mathbb{Z}_p$、$\mathbb{Q}_p^\times$，还有很多。其对偶是
\[
\widehat{G}=\Hom_{\mathrm{TopGrp}}(G,U_1),
\]
即 $G$ 的连续酉特征标。若 $G$ 有限，这恰好就是我们已经有的：所有特征标都取值于单位根。\footnote{如果你关心代数几何，也许习惯说 $\mathbb{G}_m$ 而不是 $U_1$。这个故事对 $\mathbb{G}_m$ 也有一个版本，但这种特定的正则性要求用 $U_1$。}

对我们至关重要的是：对偶 $\widehat{G}$ 是群代数的谱。在 LCA 设定中有不同的实现方式；我们的做法是给 $G$ 配一个测度。具体地说，$G$ 带有 Haar 测度，即关于 $G$ 通过乘法作用在自身上的左不变测度，且在相差一个标量倍的意义下唯一。若 $G$ 紧，这个测度是双不变的，可以归一化为总测度 $1$；但我们常常关心非紧群。

无论如何，我们用来代替 $\mathbb{C}[G]$ 的，是 $C^*$-代数 $L^1(G)$（关于选定的 Haar 测度），其卷积为
\[
(f_1*f_2)(h)=\int_G f_1(g)f_2(g^{-1}h)\,\mathrm{d}g.
\]
你可以像有限群情形那样，从核变换中提取出这个公式。有了这个定义，
$\operatorname{mSpec}(L^1(G),*)\cong\widehat{G}$，不只是作为拓扑空间，还作为群结构；这是 Fourier 变换的一个版本。这又是一种对偶，常称为 Pontrjagin 对偶，其中自然映射 $G\to\widehat{\widehat{G}}$ 是拓扑阿贝尔群的同构。

Fourier 变换仍有公式
\[
f\longmapsto\pi_{2*}(\pi_1^*f\cdot K),
\]
其中 $K$ 是核，公式与之前相同。这里的新东西是：你必须做分析，才能看出另一侧出现什么正则性。对 $L^1$，你得到同构 $(L^1(G),*)\cong(C_v(\widehat{G}),\cdot)$，其中 $C_v$ 表示在无穷远处消失的连续函数空间。Plancherel 定理给出 $(L^2(G),*)\cong(L^2(\widehat{G}),\cdot)$。而对这两者，群元素的平移都与乘法互换。（要说得完全精确，你可能得用分布，这样才能有 $\delta$-函数。）特征标与点互换；这是对称的，但在几何 Langlands 纲领以及相关的拓扑场论中，两侧并不对称。

\srcpage{15}

\begin{example}[Fourier 级数]
设 $G=U_1$。则 $\widehat{G}=\Hom(U_1,U_1)$，可以通过映射
\[
n\longmapsto(x\longmapsto\exp(2\pi i n x))
\]
与 $\mathbb{Z}$ 等同。也就是说，$\mathbb{Z}\cong\widehat{U_1}$。
\end{example}

Fourier 级数的理论把 $L^2(U_1)$ 与 $L^2(\mathbb{Z})$ 等同；后者常称为 $\ell^2$。人们把周期函数（$U_1$ 上的函数）分解成 Fourier 模。对其他类型的正则性，也有相应的版本。

要把它看成核变换，有一个 $U_1\times\mathbb{Z}$ 上的函数，取
\[
(x,n)\longmapsto\exp(2\pi i n x),
\]
而特征标 $\mathbb{Z}\to U_1$ 由 $1$ 的像决定，这个像可以取任意值，因此 $\widehat{\mathbb{Z}}=U_1$。

这照亮了 Pontrjagin 对偶的一个漂亮事实：$G$ 紧，当且仅当 $\widehat{G}$ 离散。

关于这一切，一个很好的参考文献是 Ramakrishnan–Valenza [RV99]。
